\section*{Chapitre \ref{ch:m2:fx-swaps-forwards-and-the-cross-currency-basis} --- Swaps de change, change à terme et base de devises}

\begin{solution}{exo:m2:fx-swaps-forwards-and-the-cross-currency-basis:1}
$1.1464 \times (1 + 0.0368 \times 90/360)/(1 + 0.02 \times 90/360) = 1.15119$:
$+47.9$ pips. L'euro, dont le taux est le plus bas, se traite avec une prime à terme.
\end{solution}

\begin{solution}{exo:m2:fx-swaps-forwards-and-the-cross-currency-basis:2}
Sinon il existerait un profit sans risque~: emprunter la devise à taux bas, la
convertir au comptant, prêter au taux élevé et revendre le produit à terme. Pour
que cette opération ne rapporte rien, le terme doit restituer moins de devise à taux bas par unité de
devise à taux élevé que ne le faisait le comptant~: le déport à terme de la devise à taux élevé
compense son avantage de taux.
\end{solution}

\begin{solution}{exo:m2:fx-swaps-forwards-and-the-cross-currency-basis:3}
L'investisseur a emprunté 10~millions de dollars et prêté 1\,568.70 millions de yens
pendant trois mois, chaque prêt servant de collatéral à l'autre. Il reçoit en retour 1\,557.06
millions de yens~: l'écart de 11.64 millions de yens correspond au coût des dollars, soit le taux en dollars
moins le taux en yens, auquel s'ajoute la base.
\end{solution}

\begin{solution}{exo:m2:fx-swaps-forwards-and-the-cross-currency-basis:4}
$-25$ points de base~: prêter des yens et emprunter des dollars par le swap rapporte le
taux en yens diminué de 0.25~\%. Les dollars sont plus chers
par le swap qu'au comptant, signe habituel depuis 2007.
\end{solution}

\begin{solution}{exo:m2:fx-swaps-forwards-and-the-cross-currency-basis:5}
Acheteur de dollars, le trader les vend pour l'échéance proche et les rachète pour
l'échéance suivante à 1.18 pip (0.0118 yen) de moins~: il gagne environ 1.18~million de yens par jour,
soit le carry de la détention de la devise à rendement plus élevé.
\end{solution}

\begin{solution}{exo:m2:fx-swaps-forwards-and-the-cross-currency-basis:6}
En fin de trimestre, les banques réduisent leur bilan pour leurs publications et
leurs ratios réglementaires~: la capacité d'arbitrage diminue et ceux qui ont besoin
de dollars par swaps paient davantage, si bien que la base s'élargit. En mars 2020, les lignes de swap permettent aux banques centrales
de prêter des dollars à leurs banques au taux OIS plus 25 points de base~; aucune banque disposant
d'un accès n'a eu besoin de payer beaucoup plus que cela en swaps, et la base s'est resserrée.
\end{solution}

\begin{solution}{exo:m2:fx-swaps-forwards-and-the-cross-currency-basis:7}
Décembre 2018~; 57 des 101 mois de données.
\end{solution}

\begin{solution}{exo:m2:fx-swaps-forwards-and-the-cross-currency-basis:8}
La transaction consomme du bilan~: la banque prête des dollars et détient des actifs
en yens, qui comptent dans son effet de levier et ses ratios de fonds propres, et ce coût est encore plus élevé en fin de trimestre.
Elle ne conclura l'opération que si la base rémunère ce coût~; la base est le prix du
bilan, non un profit gratuit, et la demande de couverture maintient son ouverture.
\end{solution}

\begin{solution}{pb:m2:fx-swaps-forwards-and-the-cross-currency-basis:1}
\textbf{1.} 155.8028, $-106.7$ pips.
\textbf{2.} $3.68 - 0.977 = 2.70$ points de pourcentage par an.
\textbf{3.} $4.75 - 2.70 = 2.05\%$.
\textbf{4.} 1.80~\% et 1.55~\%.
\textbf{5.} Tous inférieurs au rendement du JGB de 2.94~\%, de 0.89, 1.14 et 1.39 points.
\textbf{6.} Il clôture le swap arrivant à échéance, en payant ou en recevant des yens
pour la variation de la valeur du dollar depuis le dernier roll, puis ouvre un nouveau swap de trois mois au
nouveau comptant et aux nouveaux \omterm{def:m2:fx-swaps-forwards-and-the-cross-currency-basis:forward}{points de terme}.
\textbf{7.} Le swap gagne en yens ce que perdent les obligations~: l'assureur reçoit environ
10~\% du notionnel couvert en yens au moment du roll, compensant la baisse de la valeur en yens
de ses obligations libellées en dollars.
\textbf{8.} Lorsque le dollar monte, la couverture perd et l'assureur doit payer des yens
au moment du roll, alors que le gain sur les obligations n'est pas réalisé~; il lui faut de la trésorerie pour cela.
\textbf{9.} Il fixe le coût pour un an et évite trois roll, mais aux points et à la base
d'un an, et règle en une seule fois toute la réévaluation de l'année.
\textbf{10.} Des taux en dollars plus bas, des taux en yens plus élevés ou une base plus étroite.
\textbf{11.} 4.75~\% plus ou moins la variation du dollar, qui, une mauvaise année, peut
valoir plusieurs fois l'avantage de 1.81 point.
\textbf{12.} Le dollar ne doit pas baisser de plus d'environ 1.81~\% sur l'année face au
yen pour que le \omterm{def:m2:money-markets:bill}{bon du Trésor} non couvert égale le JGB.
\textbf{13.} Échanger une partie du rendement contre une part du risque de change, dans ses limites, et
réduire son besoin de trésorerie à renouveler.
\textbf{14.} Ce sont des demandes unidirectionnelles d'emprunt de dollars par swaps~; les
banques qui les fournissent facturent le bilan.
\textbf{15.} Les banques, et par leur intermédiaire les investisseurs disposant de dollars à prêter, comme les \omterm{def:m2:money-markets:mmf}{fonds monétaires}
prêtant en repo, et parfois les lignes de swap des banques centrales.
\textbf{16.} En octobre 2022, juste avant que les taux de la Réserve fédérale ne dépassent
suffisamment ceux du Japon pour faire passer le rendement couvert sans base sous celui du JGB.
\textbf{17.} Elle abaisse le taux en dollars et le coût de la couverture de 1 point~: le
rendement couvert monte à environ 3.05~\% sans base, au-dessus du JGB à 2.94~\%, si les
rendements obligataires ne bougent pas.
\textbf{18.} Parce qu'elle utilise le bilan de la banque, qui est coûteux~; la base
rémunère ce coût, et en fin de trimestre la banque peut ne plus prêter du tout.
\textbf{19.} Résultat à retenir~: \emph{le rendement couvert} du \omterm{def:m2:money-markets:bill}{bon du Trésor} à dix ans
pour l'assureur était d'environ 2.05~\% sans base, et de 1.80~\% avec une base de
$-25$ points de base, contre 2.94~\% pour le JGB.
\textbf{20.} Le rendement étranger diminué de l'écart de taux courts et de la base~:
la couverture reprend ce que les taux plus élevés apportent.
\end{solution}

\begin{solution}{iq:m2:fx-swaps-forwards-and-the-cross-currency-basis:1}
Emprunter une unité de la devise de base à $r_b$, la vendre au comptant à $S$, prêter
le produit à $r_q$ et convenir aujourd'hui de racheter la devise de base à $F$~: aucun coût initial,
aucun risque, donc $F(1 + r_b\tau_b) = S(1 + r_q\tau_q)$, soit
$F = S(1 + r_q\tau_q)/(1 + r_b\tau_b)$.
\iqlookfor{l'argument de réplication et les \omterm{def:m2:government-bonds:daycount}{conventions de décompte des jours}.}
\end{solution}

\begin{solution}{iq:m2:fx-swaps-forwards-and-the-cross-currency-basis:2}
Un achat au comptant de devises et sa revente à terme à une date ultérieure~: un prêt
garanti d'une devise contre une autre. Les banques financent avec lui des actifs en devises, les investisseurs
renouvellent leurs couvertures de change, les dealers gèrent leurs positions et leur liquidité, et
l'essentiel de ce marché est à très court terme, si bien qu'il est négocié encore et encore.
\iqlookfor{l'interprétation en termes de prêt et les usages.}
\end{solution}

\begin{solution}{iq:m2:fx-swaps-forwards-and-the-cross-currency-basis:3}
D'abord, en 2008 et pendant la crise de l'euro, les tensions sur le crédit et le financement des banques ont
mis fin à l'arbitrage. Depuis 2014, une forte demande de couverture unidirectionnelle pour emprunter des dollars
au moyen de swaps (investisseurs, émetteurs qui échangent des obligations en devises, banques) se heurte aux limites
de l'arbitrage~: la réglementation rend le bilan coûteux, surtout en fin de trimestre,
si bien que l'écart persiste au prix de l'occupation du bilan.
\iqlookfor{une demande et des limites à l'arbitrage, et non une cause unique.}
\end{solution}

\begin{solution}{iq:m2:fx-swaps-forwards-and-the-cross-currency-basis:4}
Vendre des euros à terme contre des dollars pour la valeur du portefeuille, en renouvelant des \omterm{def:m2:fx-swaps-forwards-and-the-cross-currency-basis:swap}{swaps de change} à un
ou trois mois et en ajustant le notionnel à mesure que la valeur évolue. La couverture
rapporte le taux du dollar moins le taux de l'euro, plus ou moins la base~: avec des taux en dollars plus élevés,
un investisseur en dollars gagne à couvrir des obligations en euros, et la base négative de l'euro
s'ajoute à ce gain. Le rééquilibrage et les flux de trésorerie à chaque renouvellement
constituent les coûts opérationnels.
\iqlookfor{le signe du carry et le rôle de la base.}
\end{solution}

\begin{solution}{iq:m2:fx-swaps-forwards-and-the-cross-currency-basis:5}
Les swaps qui couvrent la fin de trimestre deviennent plus chers~: les banques qui publient leur situation à cette date
réduisent la capacité de prêt de leur bilan, si bien que les taux implicites en dollars et la base sautent.
Une banque disposant d'une marge au bilan prête dans la hausse~; une banque qui a besoin de dollars
se finance à l'avance, avant que le \omterm{def:m2:money-markets:turn}{passage de fin de période} ne soit intégré dans les prix.
\iqlookfor{le mécanisme du bilan et une réponse par le préfinancement.}
\end{solution}

\begin{solution}{iq:m2:fx-swaps-forwards-and-the-cross-currency-basis:6}
Pour chaque devise, une courbe d'actualisation ou une courbe de marché monétaire, mise à jour à partir de ses sources~;
pour chaque paire, le comptant et les points de marché par maturité. Calculer les points de \omterm{def:m2:fx-swaps-forwards-and-the-cross-currency-basis:cip}{parité des taux d'intérêt couverte} à partir des courbes
avec les bonnes \omterm{def:m2:government-bonds:daycount}{conventions de décompte des jours} et les bonnes dates de valeur au comptant, et la base implicite
à partir des points de marché~; interpoler les dates non standard sur la base, et non sur les points~; signaler les données périmées
et les valeurs aberrantes~; publier avec horodatage~; tester sur des cas calculés à la main et
sur des allers-retours entre points et base.
\iqlookfor{conventions par devise, interpolation sur la base et contrôles.}
\end{solution}
